
This study investigated binary classification of convolutional neural network depth using weight statistics. Three feature types (global statistics, layer-wise progressions, architectural features) achieved 100\% test accuracy on 4 held-out models from a 20-model dataset spanning ResNet, VGG, and DenseNet families. Random initialization tests showed identical performance on untrained models, indicating that discriminative features reflect architectural properties (layer count, residual blocks, bottleneck ratios) rather than training-induced patterns. Within-family validation achieved 100\% accuracy for ResNet and DenseNet families.

The small test set (n=4) limits statistical confidence, and generalization to architectures beyond the tested families (Vision Transformers, modern CNNs) remains unknown. The binary classification task (shallow ≤34 vs deep ≥50) does not test continuous depth regression or fine-grained discrimination.

The results establish that structural properties create detectable signatures in weight statistics independent of training history for the tested models and task formulation. Feature importance analysis identified bottleneck ratio and layer count as primary discriminators. Global weight statistics (mean, std, min, max of layer-wise norms) achieved perfect classification despite being simple aggregations, suggesting that architectural depth is encoded in multiple weight-derived representations.

Future work could examine: (1) continuous depth regression predicting exact layer counts, (2) generalization to Vision Transformers and modern CNN architectures, (3) cross-dataset validation using models pretrained on datasets other than ImageNet, (4) larger-scale validation with 50-100 models for precise confidence intervals, (5) investigation of whether training-induced features exist beyond the architectural properties identified in this study.
